\documentclass[10pt]{article}
\usepackage[margin=20mm]{geometry}
\usepackage{amsmath,amssymb,mathtools,fontspec,unicode-math,xurl,hyperref}
\IfFontExistsTF{DejaVu Serif}{\setmainfont{DejaVu Serif}}{\setmainfont{Latin Modern Roman}}
\setmathfont{latinmodern-math.otf}
\setlength{\parindent}{0pt}\setlength{\parskip}{6pt}\sloppy
\title{Technical review of
Asymptotic stability of solitary waves for one-dimensional nonlinear Schrödinger equations}
\author{AI technical verification}
\date{}
\begin{document}
\maketitle
Abstract: The assessment of the exact acquired author version is Major. Three core proof defects are confirmed. Two essential identities are incorrect: the moving coordinate system loses a phase connection, and the negative distorted spectral branch is assigned an incorrect local cubic dynamics. Direct differentiation and a realizable mixed-scattering profile give decisive evidence. A uniform oscillatory-integral remainder used in the stationary-phase argument is also false. These are errors in the proof, not counterexamples to the main stability theorem. Consistent phase coordinates and a reduction to the positive branch give specific repair routes, whose remaining analytic obligations are identified. Several local statements also require correction, including projection hypotheses and normalization. All 106 physical source pages were read; the report distinguishes this coverage from a complete independent reconstruction of every cited theorem. Scientific adjudication remains pending.\par
\section*{Source identity and precise scope}
This review concerns M25-0269, by Charles Collot and Pierre Germain, with the title above and DOI `10.4171/jems/1642`. The acquired PDF is arXiv:2306.03668v1, submitted on 6 June 2023 at 13:31:38 UTC, with 106 physical pages and matching printed numbering. Its 936,176 bytes have SHA256\par
> 953376283a594293eb4c609cb1acf3714b08121f129b434f175f75e2f2c4cedd.\par
The 22,116-byte manifest has SHA256\par
> 061add4808e5f55f2166fd4f4fd7e4a83c3caa0fb3da1fa41008df1a5879541b.\par
The publisher landing page records submission on 22 August 2023, acceptance on 10 November 2024, and online publication on 22 April 2025. The publisher full text has not been established to equal this earlier author version. The frozen title, identifier, DOI, and 2025 batch membership are retained.\par
The equation and standing wave convention are \[iv_t-v_{xx}-F'(|v|^2)v=0,\qquad
 v=e^{-i\omega t}\Phi_\omega,\qquad
 \Phi_\omega''-\omega\Phi_\omega+F'(\Phi_\omega^2)\Phi_\omega=0.\] Theorem 1.1 assumes a smooth branch of positive even localized profiles, the positive slope \(c_\omega=\partial_\omega\int\Phi_\omega^2>0\), no discrete spectrum other than zero, and no resonance at either spectral edge. The nonlinearity has a nondegenerate minimum at zero, with \(L=F''(0)>0\). Small perturbations belong to \(H^1\cap L^{2,1}\). The conclusion includes a limiting soliton, decaying radiation, a modified scattering profile, convergence of the modulation parameters, and continuity of the scattering data. These quantifiers and the distinction between the physical radiation and its two distorted spectral components matter below.\par
\section*{Core defects}
Problem 1 (Major: omitted connection in the moving frame).  Equations (8.1)–(8.2), p. 47, and (8.9), p. 48, omit \(-y\dot p\) in the coefficient of the phase direction. This coefficient propagates to (8.26), p. 54, (9.18), p. 60, and the modulation estimates used in (12.13), Section 12.6, p. 87.\par
\subsection*{Exact differentiation and an admissible witness}
Put \(z=x+y(t)\) and \(w(t,z)=\Phi_{\omega(t)}(z)+u(t,z)\). The actual ansatz in (8.1) is \(v(t,x)=e^{i(p(t)x-\gamma(t))}w(t,x+y(t))\). Directly, \[e^{-i(px-\gamma)}(iv_t-v_{xx})
 =iw_t-w_{zz}+i(\dot y-2p)w_z+
   \bigl(\dot\gamma+p^2-\dot p(z-y)\bigr)w .\] Consequently the right-hand modulation coefficient is \[a=\omega-p^2-\dot\gamma-y\dot p,\] rather than the printed \(\omega-p^2-\dot\gamma\). The omitted vector term is \(-y\dot p(\Xi_0+\sigma_3U)\).\par
This failure occurs under the stated hypotheses of Section 8.1.1, which at this point allow arbitrary \(C^1\) frame parameters. Fix any actual soliton in the assumed branch and the exact solution \(v=e^{-i\omega t}\Phi_\omega(x)\). Near \(t_0\), choose \(p=A(t-t_0)\), \(y=b\ne0\), and \(\gamma=\omega t\), and set \[u(t,z)=e^{-ip(t)(z-b)}\Phi_\omega(z-b)-\Phi_\omega(z).\] These choices satisfy the ansatz identically. The residual of the printed scalar equation at \(t_0\) is \(-Ab(\Phi_\omega+u)\), which is nonzero for \(A\ne0\). The perturbation at \(t_0\) can be made small by taking \(b\) small. This is an exact coordinate-identity counterexample; it is not claimed to satisfy the later orthogonality choice or to furnish a counterexample to Theorem 1.1.\par
\subsection*{Impact and repair boundary}
The omitted coefficient is not automatically an integrable error. The estimates \(y=O(t)\) and \(\dot p=O(t^{-2+2\alpha})\) give only \(y\dot p=O(t^{-1+2\alpha})\). Thus the equation from which the phase estimates are derived is materially different. This affects a central derivation even though a change of phase coordinates can repair it.\par
One can consistently replace the printed phase defect by \(a\) throughout the modulation system, and then reconstruct the physical \(\gamma\). Alternatively one can change the ansatz to include the primitive of \(y\dot p\) in its phase. The geometric phase in (8.22), p. 52, can also be used, provided the conversion is explicit: if \(\gamma=-\gamma_{\rm geom}-py\), then \(a=\omega-p^2+\dot\gamma_{\rm geom}+p\dot y\). That alternative chart does not make the literal identity (8.1)–(8.2) correct.\par
There is enough conditional decay to reconstruct the desired physical phase. If the corrected equations establish \(p-p_\infty=O(t^{-1-\nu})\), \(y=O(t)\), and \(\dot y=O(1)\), then \[\int_t^\infty y(s)\dot p(s)\,\mathrm ds
 =-y(t)(p(t)-p_\infty)
   -\int_t^\infty\dot y(s)(p(s)-p_\infty)\,\mathrm ds
 =O(t^{-\nu}).\] This verifies the needed tail estimate. A fully consistent rewrite of every subsequent phase variable and of the final continuity argument is still required.\par
Problem 2 (Major: incorrect negative-branch cubic identity).  The identity (13.10), p. 92, its use in Lemma 13.5, p. 98, and the common local phase evolution for both signs in Proposition 13.1 and pp. 99–100 are false for the negative spectral branch. The reality relation (8.16), p. 50, mixes the two frequency signs through scattering.\par
\subsection*{A profile satisfying the spectral reality constraints}
For \(\xi>0\), write \[F_+=\binom{A}{B}=\binom{\widetilde f_+(\xi)}{\widetilde f_+(-\xi)},
 \quad F_-=\binom{C}{D},\quad
 R=\begin{pmatrix}r&s\\s&r\end{pmatrix},\quad
 J=\mathop{\mathrm{diag}}(1,-1).\] Here \(r=r(\xi)\), \(s=s(\xi)\), and the conventions of (5.31) and (8.16) give \[R^*R=I,\qquad F_-=-R^*\overline{F_+}.\] Expanding the actual coefficient table in (13.10) gives \[\Theta_-(\xi)=\bar r\,|A|^2\bar A+
                      \bar s\,|B|^2\bar B .
 \tag{1}\] It is not in general \(-|C|^2C\).\par
The paper’s own scattering properties imply the existence of an actual frequency \(\xi_*>0\) with \(|s(\xi_*)|^2=|r(\xi_*)|^2=1/2\): use continuity, \(s(0)=0\), \(s(\xi)\to1\), and unitarity. Thus the example does not rely on assuming that an arbitrary unitary matrix is realized by a potential. Choose a smooth compactly supported positive-branch profile near \(-\xi_*\), with value \(a\ne0\) there and zero near \(+\xi_*\). Define the negative branch by the exact reality relation. These profiles have zero trace at the origin, all required weighted and \(H^1\) regularity, and arbitrarily small norms after scaling. The inverse distorted transform is a legitimate continuous spectral physical-real vector.\par
At \(\xi_*\), \(A=0\), \(B=a\), \(C=-\bar s\bar a\). Equation (1) gives \(\bar s|a|^2\bar a\), whereas the printed expression gives \(|s|^2\bar s|a|^2\bar a\). Their difference has modulus \(|a|^3/(2\sqrt2)\). This disproves the algebraic identity on its admissible profile class. It does not assert that this chosen profile extends to an entire trapped nonlinear solution.\par
An independent exact check summed every sign term in the coefficient table. At the auxiliary algebraic fixture \(s=3/5\), \(r=4i/5\), \(A=0\), \(B=1\), the positive branches agree with the printed reduction. The two negative differences are \(48/125\) and \(-36i/125\). This rational fixture checks the algebra; the preceding intermediate-value argument supplies the actual-potential witness.\par
\subsection*{Correct dynamics and implications}
Let \(N(z_1,z_2)=(|z_1|^2z_1,|z_2|^2z_2)\) and \(k=L/(2t)\). If the verified positive equation is \[i\partial_tF_+=kN(F_+)+\tau\xi JF_++E,\] then conjugation and the reality relation force \[i\partial_tF_-=
 -kR^*N(RF_-)-\tau\xi R^*JR\,F_-+R^*\bar E .
 \tag{2}\] The leading cubic term is of order \(t^{-1}\) and cannot be discarded into an integrable remainder. Even in the reflectionless limit its sign differs from the common printed sign. For mixed scattering, two positive-branch phases are combined; there is no single local negative phase of the asserted form.\par
The linear \(\tau\) discrepancy need not invalidate the separate modulation estimate. Indeed \[R^*JR+J=
 \begin{pmatrix}2|r|^2&2\bar r s\\2\bar s r&-2|r|^2\end{pmatrix}.\] The decay and derivative bounds for \(r\) make \(\xi(R^*JR+J)\) an order-zero bounded multiplier with the needed derivative control. It can be assigned to the integrable remainder. The cubic discrepancy cannot.\par
A viable repair is to prove the pointwise bound and modified scattering only for \(F_+\), then reconstruct \(F_-\) using reality. In particular \(\|F_-\|_\infty\le\sqrt2\|F_+\|_\infty\), paired \(L^2\) norms are preserved, and derivatives of \(R\) are controlled. The first physical component in Section 14 has a propagating positive-branch contribution; the opposite branch occurs through localized eigenfunction remainders. This explains why the final physical \(L^2\) conclusion may survive. The single-phase assertion for both signs must be withdrawn, and the plus signs in the positive-branch formulas (14.1)–(14.2) must agree with the negative phase in Theorem 1.1. A complete proof of uniform profile regularity and continuity after these changes is not certified here.\par
\section*{Additional defects and independently checked boundaries}
Problem 3 (Minor: missing continuous-spectral hypotheses).  Lemma 5.15, p. 38, and Proposition 6.3, pp. 40–41, state identities in terms of the distorted transforms of unrestricted inputs. They need \(P_{\mathrm e}f=f\) and, for the cubic identity, the corresponding conditions on \(g,h\).\par
Here is an explicit operator satisfying (H1)–(H3). For \(a=1/2\), put \[L_a=-\partial_x^2-a(a+1)\mathop{\mathrm{sech}}^2x,\quad
 H=\mathop{\mathrm{diag}}(L_a+1,-L_a-1),\quad \phi=\mathop{\mathrm{sech}}^a x.\] Its potentials are even, smooth and exponentially decaying. With \(A_a=\partial_x+a\tanh x\), \[L_a=A_a^*A_a-a^2,\qquad
 A_aA_a^*-a^2=-\partial_x^2+a(1-a)\mathop{\mathrm{sech}}^2x\ge0.\] Intertwining shows that the only negative eigenvalue of \(L_a\) is \(-a^2\), with eigenfunction \(\phi\). There is no bounded zero-energy solution: for such a solution \(u\), \(A_au\) solves the positive-partner zero-energy equation. Integration by parts, with vanishing boundary derivative terms, makes both its kinetic and positive potential energies zero. Hence \(A_au=0\), which is incompatible with \(L_au=0\) unless \(u=0\). Positive-energy square-integrable solutions are excluded by the short-range scalar Jost asymptotics. It follows that \(H\) has only the two discrete eigenvalues \(\pm3/4\), no embedded eigenvalues, and no threshold resonance.\par
Take \(f=\phi e_1\), a Schwartz discrete eigenvector, and \(W(x)=xe^{-x^2}\). Then \(\widetilde f=0\), but \(Wf\) is nonzero and odd, hence entirely continuous spectral. Its distorted transform is nonzero. The two sides of Lemma 5.15 are therefore different. For Proposition 6.3 use \(f=g=h=\phi e_1\), \(j=k=\ell=m=1\), and the same \(W\), with \(L=0\). The left side is the transform of the nonzero odd vector \(W\phi^3e_1\); the right side is zero.\par
These are literal failures of the stated input scope, not failures of the transform on \(P_{\mathrm e}L^2\). The main application already separates discrete and continuous terms, notably on pp. 63 and 88. Add the projection assumptions and keep all discrete contributions in physical space. In particular, one should not apply a continuous spectral norm to a discrete basis vector as in the shorthand estimate around (9.23), p. 62. Direct localized \(H^1\) estimates of \(W\Xi_jU_{\mathrm e}\) give the required bound. No counterexample to the resulting nonlinear estimate is asserted.\par
Problem 4 (Minor: local linearization and illustrative examples).  The sign of \(F''(\Phi^2)\Phi^2\) in \(V_1\) on p. 11 is inconsistent with the actual linearization. The illustrative embedded-eigenvalue argument on p. 10 also has the wrong scalar operator. The cubic–quintic example on p. 5 uses the wrong primitive \(F\).\par
For \(F(z)=z^2\), \(\omega=1\), \(\Phi=\mathop{\mathrm{sech}}x\), the soliton equation holds and \(c_\omega=1/\sqrt\omega>0\). The printed p. 11 choice gives \(V_1=0\) and \(V_2=-2\mathop{\mathrm{sech}}^2x\), and consequently \[H(\Phi,-\Phi)^\top=4\mathop{\mathrm{sech}}^3x(1,1)^\top\ne0.\] This contradicts the claimed zero mode under the local linearization hypotheses. The correct formula is \(V_1=-F'(\Phi^2)-F''(\Phi^2)\Phi^2\). It is already used on p. 4 and in (8.7), pp. 47–48. The pure cubic profile is used here only for that local identity, not as an example satisfying the main theorem’s absence-of-resonance assumption.\par
For the p. 10 example, take \(V_1=2\mathop{\mathrm{sech}}^2x\), \(V_2=0\). The asserted scalar operator \(-\partial_x^2-V_1\) has eigenpair \((\mathop{\mathrm{sech}}x,-1)\), but \[H(0,\mathop{\mathrm{sech}}x)^\top=(0,-4\mathop{\mathrm{sech}}^3x)^\top,\] which is not an eigenvector relation. A correct embedded example is \(V_1=-6\mathop{\mathrm{sech}}^2x\), \(V_2=0\): \((-\partial_x^2+V_1)\mathop{\mathrm{sech}}^2x=-4\mathop{\mathrm{sech}}^2x\) and \(H(0,\mathop{\mathrm{sech}}^2x)^\top=3(0,\mathop{\mathrm{sech}}^2x)^\top\). Thus the existence claim remains true after replacing its example.\par
Finally \(F(z)=z^3-z^5\) produces \(3|v|^4v-5|v|^8v\) and has \(F''(0)=0\). It is not the stated cubic–quintic case. The expression \(F(z)=z^2-z^3\) used on p. 12 supplies that case. These inconsistencies are local; the later principal linearization has the correct sign.\par
Problem 5 (Minor: three defective steps in the linear estimates).  A derivative bound in (4.19), p. 18, the translation argument in Lemma 4.8, pp. 24–25, and an integration by parts in Lemma 7.3, p. 45, require correction.\par
Write \(b(\xi)=\sqrt{\xi^2+2}\), as in the paper. The asserted bound for derivatives of \(e^{-zb(\xi)+iz\xi}\) with exponential factor \(e^{-z(b(\xi)-1)}\) is false. For \(z=1\) and one derivative, the ratio of the left side to the proposed bound without its constant is \[e^{-1}b(\xi)|i-\xi/b(\xi)|\longrightarrow\infty.\] A valid bound, sufficient for every stipulated \(0<\beta<1\), is \[|\partial_\xi^k e^{-zb(\xi)+iz\xi}|
 \le C_k b(\xi)^{-k}e^{-(3/4)zb(\xi)},\qquad z\ge0.\] It follows by scaling \(zb(\xi)\) and absorbing its polynomial factors into the remaining exponential. The inequality \(3\sqrt2/4>1\) preserves the required decay margin.\par
The translation \(f_1(x,\xi)=\check f_1(x+100,\xi)\) on p. 25 does not preserve the prescribed outgoing coefficient. Already for \(V=0\) it has leading term \(e^{100i\xi}e^{ix\xi}\) instead of \(e^{ix\xi}\), contrary to (4.43). This example meets the hypotheses of that Jost construction, which does not require the later threshold assumption. Multiply by \(e^{-100i\xi}\). The translated exponentially decaying component then has the factor \(e^{-100(b(\xi)+i\xi)}\), with the derivative bounds needed there.\par
Finally the integration by parts on p. 45 must include \[-\frac1{2it}\int e^{it\xi^2}
 \left(\frac{g'F+gF'}{\xi}-\frac{gF}{\xi^2}\right)\,\mathrm d\xi .\] The last summand is omitted in the displayed expansion. After division by \(\langle x\rangle\), the eigenfunction estimate \(|F(x,\xi)|\lesssim\min(\langle x\rangle|\xi|,1)\) bounds the omitted term by \(C t^{-1}\int |g(\xi)|/|\xi|\,\mathrm d\xi\). For \(g(0)=0\), the integral on \(|\xi|<1\) is bounded by \(C\|g'\|_2\) using \(|g(\xi)|\le\sqrt{|\xi|}\|g'\|_2\), and the complementary integral by \(C\|g\|_2\). The stated local decay rate therefore survives.\par
Problem 6 (Minor: spectral normalization is inconsistent).  The kernel in Lemma 5.15, p. 38, and the scalar bilinear kernel (6.6), p. 40, have incorrect normalization. The cubic notation changes normalization between (6.10)–(6.12), (8.20), (12.4), and Section 13.1.\par
The forward and inverse distorted transforms each contribute \((2\pi)^{-1/2}\), with \(\sigma_3\) in the forward pairing and the branch sign in the inverse. For continuous inputs the multiplication kernel is exactly \[s^W_{\rho\lambda}(\xi,\eta)
 =\frac{\lambda}{2\pi}
   \langle W\psi_\lambda(\eta),\sigma_3\psi_\rho(\xi)\rangle.\] The printed kernel instead has \(\lambda/\sqrt{2\pi}\) and omits \(\sigma_3\). The scalar bilinear kernel is \[m^W_{jk,\lambda\mu}(\eta,\sigma)
 =\frac{\lambda\mu}{2\pi}
   \int W(x)\psi_{\lambda,j}(x,\eta)
                    \psi_{\mu,k}(x,\sigma)\,\mathrm dx .\] The output index, the output sign, and the extra \((2\pi)^{-1/2}\) in (6.6) do not belong to this scalar functional. For a concrete nonzero check use the operator in Problem Problem 3 (\#pr:projection), \(f=g=xe^{-x^2}e_1\) and \(W=e^{-x^2}\). Both inputs are odd and continuous spectral; the scalar integral is strictly positive, while the printed expression with output index \(1\) differs by \(1/\sqrt{2\pi}\).\par
The cubic distribution in (6.10) already contains \(1/(4\pi^2)\). Reusing that symbol with an additional \(1/(4\pi^2)\) in (8.20) and (12.4) counts the factor twice. On pp. 85 and 91–93 the calculation effectively uses the unnormalized symbol instead. A consistent correction is to reserve \(m\) for (6.10) and set \(\mu=4\pi^2m\) wherever an external factor is retained. The delta and principal-value stationary phase calculations then do yield the positive-branch coefficient \(L/(2t)\). This review does not claim that that final coefficient is wrong. The separate negative-branch defect in Problem Problem 2 (\#pr:branch) remains.\par
Problem 7 (Minor: local modulation coordinates and energy).  The stated reality of all discrete coefficients on p. 49, the scalar extraction in Section 10.1, p. 64, derived from (8.24)–(8.26), pp. 53–54, and the boost-energy identity on p. 56 need correction.\par
For the stated generalized null basis, physical reality makes \(a_0,a_3\) purely imaginary and \(a_1,a_2\) real. For example \(U=i\Xi_0=(i\Phi,-i\Phi)^\top\) satisfies \(\sigma_1\bar U=U\) and has \(a_0=i\), not a real coefficient. The estimates use their absolute values and are unchanged.\par
With the paper’s leading matrix entries \((M_0)_{01}=c_\omega\) and \((M_0)_{23}=-\int\Phi_\omega^2\), vector of defects \((a,-i\dot\omega,i(2p-\dot y),\dot p)\), and right side \(-b_j=-\langle B,\sigma_3\Xi_j\rangle\), solving the two rows gives \[\dot\omega=-i c_\omega^{-1}b_0,\qquad
 \dot p=\left(\int\Phi_\omega^2\right)^{-1}b_2 .\] The displayed extraction loses the first factor \(i\) and changes the second sign. This is distinct from the omitted phase connection in Problem Problem 1 (\#pr:frame); it does not alter the absolute-value decay bounds.\par
For the Hamiltonian convention in the paper, direct expansion of the kinetic energy gives \[H(e^{ipx}w)=H(w)+p^2M(w)+2pP(w),\qquad
 P(w)=\operatorname{Im}\int w_x\bar w\,\mathrm dx.\] The last two terms are absent from the p. 56 identity. The orthogonality to the translation mode removes the linear \(\Phi/u\) cross term in \(P(\Phi+u)\). Since \(p=O(\epsilon_1)\), the remaining boost contributions are bounded by \(C\epsilon_1^2+
C\epsilon_1\|u\|_{H^1}^2\) and can be absorbed in the coercive estimate. This verifies a local repair, without claiming that the uncorrected conservation identity is true.\par
Problem 8 (Minor: cutoffs at subunit scales and phase size).  The inhomogeneous decomposition of Section 2.2, p. 9, does not define the subunit high-frequency cutoffs used on pp. 65 and 69–70. A product of dyadic scales is also mistakenly used for a maximum on p. 79.\par
The explicit sum \(\varphi_{<\ell}=\sum_{0\le k<\ell}\varphi_k\) is zero for \(\ell<0\), and therefore \(\varphi_{>-m\kappa}=1\) for \(m>0\) under the stated partition. It cannot exclude frequencies near zero. Furthermore the precise smooth support requirements \(\mathop{\mathrm{supp}}\varphi_0=[-1,1]\) and \(\mathop{\mathrm{supp}}\varphi_1\subset[1,4]\) leave every summand zero at \(x=1\), so they cannot form a smooth partition of unity. This is an inconsistency in the definitions, not an attempt to supply a counterexample satisfying an impossible partition assumption.\par
A concrete repair uses an overlapping inhomogeneous partition for time and macroscopic frequency, and an independent rescaled cutoff for actual small frequency. Let \(\chi=1\) near \([-1,1]\), \(\chi=0\) outside \([-2,2]\), \(\delta=t^{-1/2}\), and \[f_L(\xi)=\chi(\xi/\delta)f(\xi),\qquad f_H=f-f_L.\] For \(f(0)=0\) and \(D=\|f'\|_2\), elementary integration gives \[\|f_L\|_1\lesssim\delta^{3/2}D,\quad
 \|f_L\|_2\lesssim\delta D,\quad
 \|f_H'\|_2\lesssim D,\] \[\|\xi^{-1}f_H'\|_1\lesssim\delta^{-1/2}D,\qquad
 \|\langle \xi\rangle\xi^{-1}f_H'\|_2\lesssim\delta^{-1}D .\] These prove the five intended powers in (10.7), including the derivative of the cutoff. An identical construction at scale \(t^{-\kappa}\) fixes the earlier high-frequency split. Merely changing every index to range over all integers would require new summation checks; that shortcut is not justified here.\par
On p. 79 the relevant scale is \(2^{\max(j,k_1,k_2)}t^{-\kappa}\), not \(2^{j+k_1+k_2}t^{-\kappa}\). For \(\xi=\eta=\sigma=R\) in the positive quadratic channel, the phase is \(-\omega-R^2\); a lower bound proportional to \(R^3\) is impossible as \(R\to\infty\). Writing \(M=\max(|\xi|,|\eta|,|\sigma|,1)\), the correct nonresonant bound \(|\Phi|\gtrsim Mt^{-\kappa}\) and \(|\eta\sigma|\lesssim M^2\) yield \[|\partial_\eta\partial_\sigma(\Phi^{-1})|
 \lesssim t^{3\kappa}/M .\] The analogous differentiated cutoffs have the same usable scale. These repairs address the identified local steps; they are not a claim to have rederived every omitted term in Sections 10–13.\par
Problem 9 (Major: false uniform oscillatory-integral remainder).  Lemma 13.4, pp. 96–98, has both an incorrect leading constant in its statement and a false uniform remainder. The optimization concluding Lemma 13.3, p. 96, also does not produce the exponent stated in that lemma.\par
\subsection*{The leading constant has a direct fixed-function test}
The convention is \(\widehat\phi(0)=1/\sqrt{2\pi}\). Set \(\rho=\xi=1\) and \(F=a\psi\), where \(\psi\in C_c^\infty((-1/4,1/4))\), \(\psi(0)=1\), and \(a\ne0\) is sufficiently small. Then \(F(1)=0\) and \(\|F\|_\infty+t^{-\alpha}\|F\|_{H^1}\lesssim1\) for all \(t\ge1\). On this support the substitution \(u=2p-p^2\) is monotone. It turns the principal-value integral into \(\int e^{-itu}b(u)/u\,\mathrm du\) with \(b(0)=a/\sqrt{2\pi}\). Splitting \(b\) into \(b(0)\) times a cutoff and a function divisible by \(u\) proves the limit \[-i\pi b(0)=-i\sqrt{\pi/2}\,a .\] For \(0<\alpha<1/2\) and \(0<\kappa<1\), all printed errors tend to zero, but the printed leading term is \(-ia/2\). The correct constant is already obtained in the proof on p. 97 and used on pp. 91 and 93. Thus this statement error does not change the verified positive-branch cubic coefficient.\par
\subsection*{A counterexample to the remainder even after fixing the constant}
Fix \(0<\alpha<1/4\), \(\kappa=\alpha/2\), \(\rho=1\), and a nonzero nonnegative \(\eta\in C_c^\infty((1,11/10))\). For large \(t\) put \[\xi_t=t^{-\alpha},\quad w_t=t^{-1+\alpha},\quad
 a_t=t^\alpha\sqrt{w_t}=t^{-1/2+3\alpha/2},
 \qquad F_t(p)=c\,a_t\,\eta(p/w_t).\] Here \(c>0\) is fixed and small. Both \(F_t(0)\) and \(F_t(\xi_t)\) vanish for all sufficiently large \(t\). Moreover \[\|F_t'\|_2=c\,t^\alpha\|\eta'\|_2,\qquad
 \|F_t\|_2=c\,t^{-1+2\alpha}\|\eta\|_2,
 \qquad \|F_t\|_\infty\lesssim c.\] One may set \(F_t=0\) on an initial finite time interval, with a transition after the supports are separated. Every printed hypothesis then holds for all \(t\ge1\).\par
Substitution \(p=w_tq\) and dominated convergence give \[\frac{K(t,\xi_t)}{c\,a_t}\longrightarrow
 \frac1{\sqrt{2\pi}}\int e^{-2iq}\frac{\eta(q)}q\,\mathrm dq
 \ne0.\] The real part is strictly negative because \(\cos(2q)<0\) on the support. The three printed errors have respectively the powers \[-\tfrac12+\tfrac54\alpha-\tfrac12\alpha^2,\qquad
 -\tfrac12+\alpha,\qquad
 -\tfrac34+\tfrac52\alpha .\] The actual power \(-1/2+3\alpha/2\) exceeds them by \(\alpha/4+\alpha^2/2\), \(\alpha/2\), and \(1/4-\alpha\), all strictly positive. This is a uniform-estimate counterexample satisfying all of the lemma’s hypotheses, not a numerical asymptotic guess.\par
The defective local step on p. 97 uses \(t^{-1/2}\) where \(H^1\) gives only \[|F(t^{-1/2}P)-F(0)|
 \le t^{-1/4}|P|^{1/2}\|F'\|_2 .\] Its corrected contribution is \(O(t^{\alpha-1/4})\). After the external \(t^{-1}\) in the application, this is \(O(t^{-5/4+\alpha})\), which is smaller than \(t^{-15/14}\) for \(\alpha<5/28\). For the paper’s small choice of \(\alpha\), this particular corrected contribution also fits the high-frequency error budget in Lemma 13.5. This checks only one term, however; it does not supply a complete replacement for the false uniform remainder. Proposition 13.2 explicitly invokes Lemma 13.4 on p. 94, and Lemma 13.5 invokes Proposition 13.2 on p. 98 to justify the pointwise cubic dynamics. Thus a core uniform estimate is false on its full stated hypothesis class and the required replacement has not been established here. The assessment of this component is Major, even though the leading-constant correction and the local Hölder correction are explicit. No counterexample to Proposition 13.2 on its more structured inputs, to an entire nonlinear trajectory, or to Theorem 1.1 is asserted.\par
\subsection*{A complete replacement for the insufficient optimization}
The last line of the proof of Lemma 13.3 bounds the error by \(t^{-1-\delta}+t^{-5/4+\alpha+5\delta/2}\). Balancing gives \(\delta=(2/7)(1/4-\alpha)\), not the stated \((1/3)(1/4-\alpha)\) gain. This arithmetic failure does not by itself refute the lemma. The following independent Fourier argument supplies the needed estimate.\par
After reflecting the three frequency variables, put \(P=\alpha_{\rm sign}\xi-p\) and \[H_p(t)=\iint e^{2itAB}
  g_1(P-B)g_2(B-A-P)g_3(A+P)\,\mathrm dA\,\mathrm dB.\] Write \(h_j=\mathcal F^{-1}g_j\) with the unitary convention. Gaussian regularization of the \(A,B\) integral gives exactly \[\begin{split}
 H_p(t)=\frac{\pi}{t}(2\pi)^{-3/2}
 \iiint &e^{-iP(x_1-x_2+x_3)}
 e^{-i(x_2-x_3)(x_1-x_2)/(2t)}\\
 &\hspace{8mm}\cdot h_1(x_1)h_2(x_2)h_3(x_3)
                  \,\mathrm dx_1\,\mathrm dx_2\,\mathrm dx_3 .
 \end{split}
 \tag{3}\] For example the regularized Gaussian kernel is \[\frac{\pi}{\sqrt{\epsilon^2+t^2}}
 \exp\!\left(
 -\frac{\epsilon(u^2+v^2)+2ituv}
        {4(\epsilon^2+t^2)}\right),
 \quad u=x_2-x_3,\ v=x_1-x_2.\] It converges to \((\pi/t)e^{-iuv/(2t)}\), fixing both the sign and normalization in (3).\par
Replace the second exponential in (3) by \(1\). The result is \((\pi/t)g_1(P)g_2(-P)g_3(P)\). For \(0<\delta<1/4\), use \(|e^{iz}-1|\le C_\delta|z|^\delta\) and \[|(x_2-x_3)(x_1-x_2)|^\delta
 \lesssim\prod_{j=1}^3\langle x_j\rangle^{2\delta},
 \qquad
 \|\langle x\rangle^{2\delta}h_j\|_1
 \le
 \left(\int\langle x\rangle^{4\delta-2}\,\mathrm dx\right)^{1/2}
 \|\langle x\rangle h_j\|_2 .\] The error is \(O_\delta(t^{-1-\delta+3\alpha})\), uniformly in \(P\). Taking \(\delta=1/8\) and \(\alpha<1/64\) gives a gain strictly larger than \((1/3)(1/4-\alpha)\), with difference \((1-64\alpha)/24>0\). Schwartz approximation extends the identity to the stated \(H^1\) and weighted \(L^2\) class: the latter gives \(g_1\in L^1\), and the original integral is controlled by \(\|g_1\|_1\|g_2\|_2\|g_3\|_2\). Thus this replacement is a verified repair of the identified Lemma 13.3 proof gap.\par
\section*{Coverage, unresolved obligations, and provenance}
Every physical text page 1–106 was read by the primary reviewer. The review covered the following dependencies.\par
Source pages | Actual scope and limits\par
1–12 | The theorem’s quantifiers, spectral restrictions, nonlinearity, examples, conventions, null modes, and counterexample boundaries.\par
13–30 | Jost constructions, threshold normalization, Wronskians, scattering relations, resolvent representation, and dependence on the spectral parameter. Cited external spectral theorems were not wholly reproved.\par
31–46 | Distorted inversion and weighted norm bounds, commutators, multiplication, quadratic and cubic distributions, dispersive and local decay estimates.\par
47–64 | Moving coordinates, discrete/continuous splitting, reality, orthogonality, modulation matrix, conservation bounds, bootstrap, and separation of localized errors.\par
65–88 | Low-frequency and nonresonant decompositions, modulation integration, quadratic normal forms, weighted estimates for cubic and higher terms, and exponent budgets. Identified false cutoffs and local scale bounds were separated from unproved global replacements.\par
89–100 | Stationary phase, delta and principal-value terms, the full cubic branch algebra, and pointwise evolution. The exact sign table and an independent Fourier repair were reconstructed.\par
101–106 | Physical modified scattering, modulation tails, continuity assertions, appendix, and bibliography. Some stronger regularity conclusions remain insufficiently verified, as described below.\par
The primary reviewer actually viewed 67 source evidence images, including every page used for the confirmed issues. Their individual byte counts and hashes are in the source inspection record. Rendering receipts remain separate from later actual-view records. Full text reading does not mean that every estimate or every external reference has been independently proved.\par
In particular, the inference near (14.7), p. 101, from an \(o(L^2)\) statement and Plancherel alone does not supply an \(o(H^1)\) conclusion. The physical scattering claim there uses \(L^2\), but the stronger profile regularity and continuity in Theorem 1.1 require their own uniform high-frequency argument. No actual nonlinear counterexample to those assertions has been established, so this is recorded as an unresolved verification obligation rather than a confirmed false theorem. The passage involving the limiting time-dependent linear phase in Proposition 13.1 also has not been fully reconstructed in the strongest stated topology. The cited Jost/resolvent theory and the admissibility of the cubic–quintic examples remain partly external inputs.\par
Three existing native AI auxiliaries performed bounded checks in this same task. Their actual text/image coverage, calculations, unsuccessful retrievals, and limitations are preserved individually. The primary reviewer read their final records and independently checked the decisive equations and evidence pages. These are AI assistance records, not human reviews or independent acceptance. Some historical auxiliary strings lost non-ASCII characters in serialization; those records are preserved verbatim and the canonical identity and mathematics here were checked against the source.\par
The primary independent SymPy 1.13.3 calculation checks the exact coordinate differentiation, the scalar eigenfunction identities, all 64 cubic sign summands, the scattering-matrix discrepancy, and the exponent comparisons. The first execution exited with code 1 because a time-dependent substitution was incorrectly inserted into a formal partial derivative in the checking script. The corrected second script exited with code 0. Both scripts, outputs, errors, and actual exit records are retained. No author code was executed.\par
The official arXiv and EMS records were checked and targeted searches for a correction, supplement, and code were made. No applicable additional item was identified in those bounded searches. This is not proof of nonexistence, and absent code is not a defect of this mathematical paper. No publisher-version equivalence is presumed.\par
The dispatch token is\par
The exact source, manifest, current TeX, compiler, dependencies, compiled PDF, images, logs, and scientific records are bound in the supporting inventory. The external controller supplied the UI observations; this executor did not re-observe the account or model selector. The runtime identifies itself as ChatGPT GPT-6 Astra Pro, but backend equivalence to the reported UI is not independently verified. The current scientific decision is `pending`; `human\_reviews=[]`. The assessment Major records the three core proof defects above. Current independent acceptance is recorded in the associated provenance; neither publication approval nor a counterexample to the main theorem is asserted.\par
9 C. Collot and P. Germain, Asymptotic stability of solitary waves for one-dimensional nonlinear Schrödinger equations, arXiv:2306.03668v1 (6 June 2023), exact acquired source. https://arxiv.org/abs/2306.03668. EMS Press, publisher metadata for DOI 10.4171/jems/1642, online publication 22 April 2025. https://ems.press/journals/jems/articles/14298719.\par
\end{document}
